\section{Results}
\label{sec:results}

\begin{figure}[t]
\centering
\includegraphics[width=\columnwidth]{decisions}
\caption{The share of clear decisions that each miniature keeps, with the count and the cost per
run. A cost of 1$\times$ is the cost of the full workflow. The dashed line is the noise ceiling: the full
workflow, run a second time. W2 has no ceiling. Section~\ref{sec:ceiling} discusses the writer group
of W3.}
\label{fig:decisions}
\end{figure}

\begin{figure*}[t]
\centering
\includegraphics[width=\textwidth]{tracking}
\caption{Each dot is one variant: its mean score on the full workflow ($x$) and on the miniature
($y$). Hollow dots are mutants. A valid miniature keeps the order of the dots (a, b). An invalid
miniature gives all variants about the same score (c, d). $r$ is the Pearson correlation.}
\label{fig:tracking}
\end{figure*}

Figure~\ref{fig:decisions} gives the main result for each miniature. Figure~\ref{fig:tracking}
shows why some miniatures keep decisions and others do not. Appendix~\ref{app:all} gives all the
measures for each miniature.

\subsection{Simulated tools keep decisions}
Simulated tools were the best shortcut of group~C. For the Wikipedia agent (W1), they kept 19 of 20
clear decisions when grounded in recorded outputs, and 18 of 20 when pure. The noise ceiling kept
20 of 20. For the research planner (W3), simulated researcher tools kept all 8 clear decisions,
and the second full run kept only 6 of 8. In Figure~\ref{fig:tracking}a--b, the scores of the
miniature follow the scores of the full workflow ($r = 0.95$ and $0.74$).

The grounded simulator also kept the diagnosis better. The failure reasons of the judge agreed
with the full workflow at $\pi = 0.41$ for the grounded simulator and at $0.27$ for the pure
simulator. The ceiling is $0.57$. A pure simulator invents pages that do not contain the facts, so
it shows more retrieval failures. Thus, a simulator must use real recorded outputs when it can.

\subsection{Replay and reduction are cheap and safe}
The cheapest valid miniatures came from groups~A and~B. A random quarter of the tasks kept 17 of 20
clear decisions at 0.20$\times$ the cost, and it combines with any other shortcut. Short tool
outputs kept all 20 clear decisions. A cheaper simulated customer in W2 kept 11 of 14, with a
ranking agreement of $\tau = 0.73$. Replay of the planner and the researchers in W3 cost
0.21$\times$ the full run and gave the same outputs for the same inputs.

\subsection{The noise ceiling protects a good miniature}
\label{sec:ceiling}
The writer group of W3 shows why a team must measure the noise ceiling. The replay miniature gave
each writer variant the same research. Its ranking of the writer variants disagreed with the full
workflow ($\tau = -0.47$). But the full workflow also disagreed with its own second run
($\tau = -0.05$), and only 2 of its 5 clear decisions came back. Each full run collects new
research, and this noise is larger than the effect of the writer variants.

Two runs of the replay miniature differed by only 0.005 in mean score. Under replay, all writer
variants scored 0.82--0.85. This agrees with the full workflow, where a mutant that ignores the
research notes scored 0.81 and the base scored 0.82. Thus, the replay miniature was correct: the
writer variants differ very little. Without the ceiling, we can reject this good miniature by
mistake.

\subsection{Two limits}
\label{sec:limits-found}
Two shortcuts failed. Both failures have one cause: the miniature must imagine information that
only the real system has. These limits are useful, because they tell a team where not to use a
shortcut.

\paragraph{Merged calls.}
One model call wrote the full agent run and imagined each tool result. This miniature kept 10 of
20 clear decisions, the same as chance. In Figure~\ref{fig:tracking}c, all variants get about the
same score. A limit of two agent steps drops the real agent from 0.74 to 0.40. The merged call
still scores 0.62, because the model imagines the facts that it did not have the steps to find. In
W3, stubbed researchers kept 0 of 8 clear decisions and a full pipeline in one call kept 4 of 8.

\paragraph{Private data.}
In W2, the tools read and write a private shop database. A simulator cannot know which customers
and orders exist. Box~1 shows a real example. In 260 of 320 runs with simulated tools, the agent
did not find the customer and sent the customer to a human agent. With real tools, this
occurred in 63 of 320 runs. All variants then scored 0.07, so the miniature kept no clear decisions
(Figure~\ref{fig:tracking}d).

\begin{note}
\textbf{Box 1. W2, task \texttt{retail-10}.} The customer wants to return all items.\\[2pt]
\textcolor{simgreen!70!black}{\textbf{Real tool.}}
A search by the e-mail address of the customer $\rightarrow$ \texttt{mia\_garcia\_4516}.
The agent finds the two orders of the customer and starts the return.\\[2pt]
\textcolor{mergeorange}{\textbf{Simulated tool.}} The same call $\rightarrow$
\texttt{"Email not found"}. The agent asks for a zip code. The customer cannot give it, and the
agent sends the customer to a human agent.
\end{note}

\subsection{Diagnoses and fixes carry back}
A miniature is more useful if a team can also learn from its failures. A blind LLM instance read
the failures from three sources and wrote three fixes from each source. Fixes from the full
workflow improved it by 2.9 points on average. Fixes from the grounded simulator improved it by
1.3 points, and fixes from the merged calls changed it by $-0.8$ points. This order follows the
reason agreement, but no single fix moved the score by more than two standard errors. We take this
result as first evidence, not as proof.

\paragraph{Another model family.}
We also read a run of W1 on a model of another family as a miniature of the main run. It kept 13
of 14 clear decisions. Its failure reasons were different: the most frequent reason was the same
for only 15\% of the variants. Thus, a cheaper model can choose between variants, but it can give
different diagnoses.
